\documentclass[11pt,reqno]{amsart}

\usepackage{amsmath,amssymb,amsthm,amsfonts,mathtools}
\usepackage[hidelinks]{hyperref}
\usepackage[a4paper,margin=1in]{geometry}

% Tolerate long, unbreakable \texttt{...} identifiers without
% producing visible right-margin overruns.
\tolerance=2000
\emergencystretch=5em
\hbadness=2000
\hyphenpenalty=50

% Same theorem setup as the main paper, so that all statement numbers
% (Theorem/Lemma/Corollary/Proposition/Definition/Remark/Example) coincide
% with those of the main paper.
\theoremstyle{plain}
\newtheorem{theorem}{Theorem}[section]
\newtheorem{lemma}[theorem]{Lemma}
\newtheorem{corollary}[theorem]{Corollary}
\newtheorem{proposition}[theorem]{Proposition}

\theoremstyle{definition}
\newtheorem{definition}[theorem]{Definition}

\theoremstyle{remark}
\newtheorem{remark}[theorem]{Remark}
\newtheorem{example}[theorem]{Example}

\title[Statement Compendium: Rational Cut-and-Project Period Formulas]{%
  Statement Compendium for\\
  \emph{Period Formulas for Rational Cut-and-Project Strip Projections:
  Multiset and Set Cases}}

\author{Dirk Kunert}
\address{Independent Researcher}
\email{dirk.kunert@gmail.com}

\subjclass[2020]{Primary 52C23; Secondary 11A07, 11B25, 11H06, 37B10, 68V20}
\keywords{cut-and-project sets, projected multisets, gap sequences,
period length, set-valued period, rational slope}

\date{June 13, 2026}

\begin{document}

\maketitle

\begin{abstract}
This is a companion to the paper \emph{Period Formulas for Rational
Cut-and-Project Strip Projections: Multiset and Set Cases} (D.\ Kunert).
It collects, for quick reference, \emph{every} formal statement of that
paper---definitions, theorems, propositions, lemmas, corollaries, remarks,
and examples---reproduced \emph{verbatim} and \emph{without proofs}.
Section headings and statement numbers coincide exactly with those of the
main paper, so any item here can be located there by its number.
\end{abstract}

\noindent\textbf{Conventions for this compendium.}
\begin{itemize}
  \item Statements are quoted verbatim from the main paper; proofs, figures,
        and connecting prose are omitted.
  \item Theorem/Lemma/Corollary/Proposition/Definition/Remark/Example
        numbers (e.g.\ \emph{Theorem 3.1}, \emph{Lemma 4.4}) match the main
        paper exactly. Section numbering therefore starts at~\S2; the main
        paper's Introduction (\S1) and its proof, formalization, literature,
        and discussion sections contain no formal statements and are not
        reproduced. (Equation numbers are local to this document.)
  \item Three references that point into the main paper's proofs, figures, or
        bibliography have been redirected to read ``the main paper''; every
        other cross-reference resolves within this document.
\end{itemize}

\medskip

% =====================================================================
% Section numbers start at 2 so that statement numbers match the paper.
\setcounter{section}{1}
\section{Construction of the Sequences}\label{sec:setup}
% =====================================================================

This section fixes the minimal notation needed to read the statements; the
full construction is given in the main paper. Fix coprime
\(\alpha, \beta \in \mathbb{N}\) and a window parameter \(\omega \ge 0\), and
project the accepted lattice points of the rational strip onto the line of
slope \(\alpha/\beta\). Sorting the projected values \(\tilde{p}\) in
nondecreasing order, \emph{with multiplicities retained}, yields
\((\tilde{p}^{(i)}_s)_{i \in \mathbb{Z}}\) and the \emph{gap sequence}
\begin{equation}\label{eq:delta}
    \delta^{(i)} \;:=\; \tilde{p}^{(i+1)}_s - \tilde{p}^{(i)}_s.
\end{equation}

\begin{definition}\label{def:period}
The sequence \((\delta^{(i)})_{i \in \mathbb{Z}}\) has \emph{period length}
\(\lambda \in \mathbb{N}\) if \(\lambda\) is the smallest positive integer
such that \(\delta^{(i+\lambda)} = \delta^{(i)}\) for all~\(i\).
\end{definition}

% =====================================================================
\section{Main Result}\label{sec:mainresult}
% =====================================================================

\begin{theorem}[Period length formula]\label{thm:main}
Let \(\alpha, \beta \in \mathbb{N}\) with \(\gcd(\alpha, \beta) = 1\) and let
\(\omega \in \mathbb{R}_{\ge 0}\).  Define
\begin{equation}\label{eq:ND}
    N \;:=\; \lfloor\omega\alpha\rfloor + \lfloor\omega\beta\rfloor + 1,
    \qquad
    D \;:=\; \alpha^2 + \beta^2.
\end{equation}
Then the period length of the gap sequence~\eqref{eq:delta} is
\begin{equation}\label{eq:formula}
    \lambda_{\mathrm{multiset}} \;=\;
    \begin{cases}
        N   & \text{if } D \nmid N, \\
        N/D & \text{if } D \mid N.
    \end{cases}
\end{equation}
Here \(\lambda_{\mathrm{multiset}}\) denotes the minimal period
\(\lambda\) of Definition~\ref{def:period} applied to the multiset gap
sequence; the set-valued analogue \(\lambda_{\mathrm{set}}\) is introduced
in Section~\ref{sec:setvalued}.
\end{theorem}

\begin{remark}[Interpretation of \(N\)]\label{rem:N}
The quantity \(N\) equals
\[
    \#\bigl([-\omega\beta,\,\omega\alpha]\cap\mathbb{Z}\bigr),
\]
the number of integers in the range of the invariant
\(r = \alpha x - \beta y\) (see the generic-case proof in the main paper).
Geometrically, \(N\) counts the number of tracks, namely the parallel arithmetic progressions of
projected values counted with multiplicity (illustrated in the main paper).
\end{remark}

\begin{remark}[Interpretation of \(D\)]\label{rem:D}
The quantity \(D = \alpha^2+\beta^2\) is the common difference of each
arithmetic progression of \(\tilde{p}\)-values.
It also equals \(\lVert(\beta,\alpha)\rVert^2\), the squared length of the
direction vector of~\(f\).
\end{remark}

\begin{remark}[Computational verification]\label{rem:verification}
The multiset and set formulas have been verified jointly against the same
CSV datasets, since each row records both \(\lambda_{\mathrm{multiset}}\)
and \(\lambda_{\mathrm{set}}\):
\begin{itemize}
    \item 51{,}012 test cases with \(\alpha, \beta\) up to \({\sim}100{,}000\)
          and rational \(\omega\) up to \({\sim}10\), all satisfying
          \(D \nmid N\) (51{,}011 with \(N < D\), one with \(N \ge D\));
          see file \texttt{multiset\_\allowbreak and\_\allowbreak set\_\allowbreak new\_\allowbreak find\_\allowbreak patterns\_\allowbreak 51012\_\allowbreak lines.csv}
          in the repository's \texttt{tests/} folder.
    \item 477 cases at \((\alpha, \beta) = (1, 2)\) with rational \(\omega\)
          chosen so that \(D \mid N\) (with \(N\) up to \({\sim}4{,}300\)),
          verifying the \(\lambda = N/D\) branch; see file
          \texttt{multiset\_\allowbreak and\_\allowbreak set\_\allowbreak degenerate\_\allowbreak patterns\_\allowbreak 5000\_\allowbreak lines.csv}
          in the repository's \texttt{tests/} folder.
    \item 240 additional cases generated by an independent Python
          implementation, balanced 120 / 120 between \(N < D\) and
          \(N \ge D\) and spanning 25 distinct \((\alpha, \beta)\) pairs;
          see file \texttt{set\_\allowbreak theorem\_\allowbreak balanced\_\allowbreak test.csv}.
\end{itemize}
All 51{,}729 cases produce exact matches with zero failures for both the
multiset formula~\eqref{eq:formula} and the set formula~\eqref{eq:setformula}
(see Theorem~\ref{thm:setperiod} below).  These computations are used only
as consistency checks; the proofs below are independent of the empirical
verification.

Historically, the \(D \nmid N\) branch of the multiset
formula~\eqref{eq:formula} was first suggested by inspecting the datasets
at \(\omega = 1\), where
\(\lambda_{\mathrm{multiset}} = \alpha + \beta + 1\)
(see Corollary~\ref{cor:omega1}).
\end{remark}

% =====================================================================
\section{\texorpdfstring{Proof: The Generic Case \(D \nmid N\)}{Proof: The Generic Case D does not divide N}}\label{sec:proof}
% =====================================================================

\emph{In the main paper the statements below constitute the proof of the
generic case; their own proofs are omitted here.}

\begin{lemma}\label{lem:gcd}
\(\gcd(\alpha, D) = \gcd(\beta, D) = 1\).
\end{lemma}

\begin{corollary}\label{cor:inverse}
The modular inverse \(\beta^{-1} \pmod{D}\) exists, and the map
\(r \mapsto c_r = -\alpha\beta^{-1}r \bmod D\) is a bijection on
\(\mathbb{Z}/D\mathbb{Z}\).
\end{corollary}

\begin{lemma}[Nonuniform residue distribution]\label{lem:nonuniform}
If \(D \nmid N\), the multiset of \(N\) residues
\(\{c_r \bmod D : r \in \mathcal{R}\}\) is not uniformly distributed over
\(\mathbb{Z}/D\mathbb{Z}\).  More precisely, writing \(N = qD + s\) with
\(0 < s < D\), exactly \(s\) residue classes are hit \(q+1\) times and
\(D - s\) residue classes are hit \(q\) times.
\end{lemma}

\begin{lemma}[Uniform residue distribution]\label{lem:uniform}
If \(D \mid N\), then every residue class \(c \in \mathbb{Z}/D\mathbb{Z}\) is
hit exactly \(N/D\) times by the multiset
\(\{c_r \bmod D : r \in \mathcal{R}\}\).
\end{lemma}

\begin{lemma}[Minimal period divides every period]
\label{lem:minperiod_divides_periods}
Let \((a_i)_{i\in\mathbb Z}\) be a bi-infinite sequence with minimal
positive period~\(\lambda\).  If \(M > 0\) is also a period, then
\(\lambda \mid M\).
\end{lemma}

\begin{lemma}[Trivial stabilizer of a cyclic interval]
\label{lem:cyclic_interval_stabilizer}
Let \(0 < s < D\) and let
\[
    H = \{x_0,\, x_0+1,\, \ldots,\, x_0+s-1\} \pmod{D}
\]
be a cyclic interval of length \(s\) in \(\mathbb{Z}/D\mathbb{Z}\).
If \(\tau \in \mathbb{Z}/D\mathbb{Z}\) satisfies \(H + \tau = H\),
then \(\tau \equiv 0 \pmod{D}\).
\end{lemma}

\begin{proposition}\label{prop:generic_period}
If \(D \nmid N\), then the multiset gap sequence has minimal
period \(\lambda = N\).
\end{proposition}

% =====================================================================
\section{\texorpdfstring{The Degenerate Case \(D \mid N\)}{The Degenerate Case D divides N}}\label{sec:degenerate}
% =====================================================================

\begin{theorem}\label{thm:degenerate}
If \(D \mid N\), then \(\lambda = N/D\).
\end{theorem}

\begin{example}\label{ex:degenerate}
For \(\alpha = 2\), \(\beta = 1\), \(\omega = 3\):
\[
    N = \lfloor 6 \rfloor + \lfloor 3 \rfloor + 1 = 10,
    \qquad
    D = 5.
\]
Since \(5 \mid 10\), the formula gives \(\lambda = 10/5 = 2\).
Direct computation confirms the gap sequence
\((0, 1, 0, 1, \ldots)\) with period~\(2\).
\end{example}

\begin{example}\label{ex:degenerate2}
For \(\alpha = 3\), \(\beta = 2\), \(\omega = 5\):
\[
    N = \lfloor 15 \rfloor + \lfloor 10 \rfloor + 1 = 26,
    \qquad
    D = 13.
\]
Since \(13 \mid 26\), the formula gives \(\lambda = 26/13 = 2\).
\end{example}

\begin{remark}\label{rem:boundary}
The boundary case \(N = D\), i.e. \(N/D = 1\), gives
\(\lambda_{\mathrm{multiset}} = 1\): every integer is a projected value in
the \(\tilde{p}\)-coordinate with multiplicity~\(1\), and all gaps
equal~\(1\).  At this boundary the multiset and set conventions agree
(Theorem~\ref{thm:setperiod}, case \(N \ge D\); see also
Corollary~\ref{cor:setlemultiset}).
\end{remark}

% =====================================================================
\section{The Set-Valued Period}\label{sec:setvalued}
% =====================================================================

Sorting instead the \emph{set} of distinct projected values
\(\tilde{p} \in \mathbb{Z}\), with each value listed once, yields a strictly
increasing enumeration \((\tilde{p}^{(i)}_{\mathrm{set}})_{i \in \mathbb{Z}}\)
and the \emph{set-valued gap sequence}
\begin{equation}\label{eq:delta_set}
    \delta^{(i)}_{\mathrm{set}}
    \;:=\; \tilde{p}^{(i+1)}_{\mathrm{set}} - \tilde{p}^{(i)}_{\mathrm{set}}.
\end{equation}
Its minimal period \(\lambda_{\mathrm{set}}\) is defined exactly as in
Definition~\ref{def:period}.

\begin{theorem}[Set-valued period]\label{thm:setperiod}
Let \(\alpha, \beta \in \mathbb{N}\) with \(\gcd(\alpha,\beta) = 1\) and
\(\omega \in \mathbb{R}_{\ge 0}\).  With \(N\) and \(D\) as
in~\eqref{eq:ND}, the set-valued gap sequence~\eqref{eq:delta_set}
has minimal period
\begin{equation}\label{eq:setformula}
    \lambda_{\mathrm{set}} \;=\;
    \begin{cases}
        N & \text{if } N < D, \\
        1 & \text{if } N \ge D.
    \end{cases}
\end{equation}
\end{theorem}

\begin{corollary}\label{cor:setlemultiset}
For all \(\alpha, \beta, \omega\) as above,
\[
    \lambda_{\mathrm{set}} \;\mid\; \lambda_{\mathrm{multiset}},
\]
and in particular
\[
    \lambda_{\mathrm{set}} \;\le\; \lambda_{\mathrm{multiset}},
\]
with equality if and only if \(N \le D\).
\end{corollary}

\begin{example}[\(N < D\), set and multiset agree]\label{ex:set_lt}
For \(\alpha = 2\), \(\beta = 3\), \(\omega = 2\): \(N = 4 + 6 + 1 = 11\)
and \(D = 13\), so \(N < D\).  Then
\(\lambda_{\mathrm{multiset}} = \lambda_{\mathrm{set}} = 11\).
\end{example}

\begin{example}[\(N = D\), boundary]\label{ex:set_eq}
For \(\alpha = 2\), \(\beta = 1\), \(\omega = 3/2\): \(N = 5 = D\), so
\(\lambda_{\mathrm{multiset}} = N/D = 1 = \lambda_{\mathrm{set}}\)
(cf.\ Remark~\ref{rem:boundary}).
\end{example}

\begin{example}[\(N > D\), \(D \nmid N\), set and multiset disagree]
\label{ex:set_gt}
Take \(\alpha = 3\), \(\beta = 4\), \(\omega = 4\).  Then
\[
    N = \lfloor 12 \rfloor + \lfloor 16 \rfloor + 1 = 29,
    \qquad
    D = 9 + 16 = 25,
    \qquad
    N = 1 \cdot D + 4,
\]
so the residue distribution
(Lemma~\ref{lem:nonuniform}) hits exactly \(4\) classes with
multiplicity \(2\) and the remaining \(21\) classes with
multiplicity \(1\).

\medskip
\noindent\emph{Multiset side.}  Within each length-\(D\) window of the
\(\tilde{p}\)-axis, the \(N = 29\) projected entries (counted with
multiplicity) include four duplicated values: each of the four
heavy residue classes contributes a coincident pair, hence a zero
gap.  The block of \(N\) gaps therefore contains four \(0\)'s and
\(N - 4 = 25\) positive values, summing to \(D\).
Theorem~\ref{thm:main} gives
\(\lambda_{\mathrm{multiset}} = N = 29\), and the heavy classes form
a cyclic interval of length \(s = 4\) whose only translational
symmetry is the identity
(Lemma~\ref{lem:cyclic_interval_stabilizer}).

\medskip
\noindent\emph{Set side.}  Since every residue class is hit at least
once, the underlying point set is all of \(\mathbb{Z}\) in the
\(\tilde{p}\)-coordinate.  Discarding multiplicities collapses the four
zero gaps and leaves the constant gap sequence \(1, 1, 1, \ldots\);
Theorem~\ref{thm:setperiod} gives
\(\lambda_{\mathrm{set}} = 1\).  Direct computation confirms both
values.
\end{example}

% =====================================================================
\section{Corollaries}\label{sec:corollaries}
% =====================================================================

Throughout this section, \(\lambda\) denotes the multiset period
\(\lambda_{\mathrm{multiset}}\) of Theorem~\ref{thm:main}; the
set-valued analogues follow from Theorem~\ref{thm:setperiod} via
Corollary~\ref{cor:setlemultiset}.  Define
\(\Lambda_{\alpha,\beta} := \alpha + \beta + 1\).

\begin{corollary}[Finiteness]\label{cor:finite}
For all \(\alpha, \beta \in \mathbb{N}\) with \(\gcd(\alpha,\beta)=1\) and
all \(\omega \in \mathbb{R}_{\ge 0}\), the period length~\(\lambda\) is
finite.
\end{corollary}

\begin{corollary}[Period at \(\omega = 1\)]\label{cor:omega1}
For \(\omega = 1\), \(\lambda = \Lambda_{\alpha,\beta}\).
\end{corollary}

\begin{corollary}[\(0 < \omega < 1\)]\label{cor:small_omega}
If \(0 < \omega < 1\), then
\[
    \lambda < \Lambda_{\alpha,\beta}.
\]
\end{corollary}

\begin{corollary}[\(1 < \omega < 2\)]\label{cor:medium_omega}
If \(1 < \omega < 2\) and \(D \nmid N\), then
\[
    \lambda \ge \Lambda_{\alpha,\beta}.
\]
\end{corollary}

\begin{remark}
The condition \(D \nmid N\) in Corollary~\ref{cor:medium_omega} cannot be
dropped.  For \(\alpha = 2\), \(\beta = 1\), \(\omega = 3/2\), one has
\[
    N = \lfloor 3 \rfloor + \lfloor 3/2 \rfloor + 1 = 5 = D.
\]
Since \(D \mid N\), we get \(\lambda = N/D = 1\), which is less than
\(\Lambda_{2,1} = 4\).
\end{remark}

\begin{corollary}[\(\omega > 2\)]\label{cor:large_omega}
If \(\omega > 2\) and \(D \nmid N\), then
\[
    \lambda > \Lambda_{\alpha,\beta}.
\]
\end{corollary}

\begin{remark}
In both corollaries above, the degenerate case \(D \mid N\) can produce
small periods: \(\lambda = N/D\) may be much less than
\(\Lambda_{\alpha,\beta}\).  This occurs when the window parameter is tuned so
that the number of tracks is an exact multiple of \(D = \alpha^2+\beta^2\).
\end{remark}

% =====================================================================
\section{The Irrational Case}\label{sec:irrational}
% =====================================================================

\begin{proposition}[Aperiodicity for irrational slope]\label{prop:irrational}
Let \(a \in \mathbb{R}_{>0} \setminus \mathbb{Q}\) and \(\omega > 0\), and
consider the analogous strip projection with line \(f(x) = a x\) and strip
\(\{(x,y) : a(x-\omega) \le y \le a x + \omega\}\).  Then the projected
gap sequence has \emph{no} finite period.  Moreover, the orthogonal
projection is injective on \(\mathbb{Z}^2\), so the multiset and set
conventions coincide in this case.
\end{proposition}

\begin{remark}[Period dichotomy]\label{rem:dichotomy}
Within the family of positive slopes considered here, and for \(\omega > 0\),
Corollary~\ref{cor:finite} and Proposition~\ref{prop:irrational} together
form a complete dichotomy:
\[
    \text{the gap sequence has a finite period}
    \;\;\Longleftrightarrow\;\;
    \text{the slope is rational.}
\]
There are no exceptional irrational slopes that produce a finite period;
the boundary between the periodic and the aperiodic regime coincides exactly
with the boundary \(\mathbb{Q}_{>0} \subset \mathbb{R}_{>0}\).  The
aperiodic side is the classical setting of model sets and aperiodic order
(see the related-literature section of the main paper).
\end{remark}

\end{document}
